\documentclass[11pt,a4paper]{article}

% ===== XeLaTeX + Korean =====
\usepackage{kotex}
\usepackage{fontspec}
\setmainfont{Times New Roman}
\setmainhangulfont{AppleMyungjo}
\setsanshangulfont{Apple SD Gothic Neo}

% ===== Layout =====
\usepackage[a4paper,top=18mm,bottom=18mm,left=20mm,right=20mm]{geometry}
\usepackage{fancyhdr}
\setlength{\headheight}{24pt}
\usepackage{enumitem}
\usepackage{mdframed}
\usepackage{array}

\setlength{\parindent}{1em}
\linespread{1.18}

% ===== Header / Footer =====
\pagestyle{fancy}
\fancyhf{}
\renewcommand{\headrulewidth}{0.6pt}
\fancyhead[L]{\sffamily\bfseries 2026 고1 영어 변형 — 지문 34}
\fancyhead[R]{\sffamily 결과에 따른 행동 조정}
\renewcommand{\footrulewidth}{0pt}
\fancyfoot[C]{\framebox[16mm][c]{\thepage}}

% ===== helpers =====
\newcommand{\qno}[1]{\par\vspace{8pt}\noindent{\sffamily\bfseries #1.}\ }
\newcommand{\instr}[1]{{\sffamily #1}\par\vspace{3pt}}
\newlist{choices}{enumerate}{1}
\setlist[choices]{label=,leftmargin=1.5em,itemsep=2pt,topsep=3pt,parsep=0pt}

\newmdenv[linewidth=0.6pt,roundcorner=3pt,innertopmargin=7pt,innerbottommargin=7pt,
          innerleftmargin=9pt,innerrightmargin=9pt,skipabove=5pt,skipbelow=5pt]{pbox}
\newcommand{\nemo}[1]{\fbox{\,#1\,}}

\begin{document}

\begin{center}
{\fontsize{15}{17}\selectfont\sffamily\bfseries [지문 34] 다음 글을 읽고 물음에 답하시오.}
\end{center}
\vspace{4pt}

% ===== 공통 지문 =====
\begin{pbox}
\hspace{1em}A particularly powerful way in which human beings adapt and adjust to their circumstances is by changing behaviors in light of the consequences that they yield. For instance, a child who gets praised by her teacher for doing an exercise correctly is likely to continue doing the exercise in that way. A dancer who feels that a particular conditioning exercise is not giving any results is likely to give that exercise up in favor of another one that seems more promising. In general, we tend to do things that bring us some kind of \rule{2.5cm}{0.4pt} or act in ways that help us avoid some kind of problem, discomfort, or disadvantage. We also tend to do fewer things (or stop doing things) that lead to problems, discomfort, or disadvantages, and we tend to avoid acting in ways that do not bring any benefit or reward.
\end{pbox}

% ===================== Q1 =====================
\qno{1}\instr{다음 빈칸에 들어갈 말로 가장 적절한 것은? [3점]}
\begin{choices}
\item ① difficulty
\item ② uncertainty
\item ③ advantage
\item ④ obligation
\item ⑤ distraction
\end{choices}

% ===================== Q2 =====================
\qno{2}\instr{윗글의 주제로 가장 적절한 것은?}
\begin{choices}
\item ① how the consequences of behavior shape what people choose to do
\item ② why children need external praise to develop good habits
\item ③ the difference between motivation and discipline in human behavior
\item ④ how discomfort can sometimes lead to greater personal achievement
\item ⑤ the role of repetition in building effective exercise routines
\end{choices}

% ===================== Q3 (서술형) =====================
\qno{3}\instr{[서술형] 다음 문장을 우리말로 해석하시오.}
\vspace{4pt}\par\noindent
\hspace{1em}\textit{A dancer who feels that a particular conditioning exercise is not giving any results is likely to give that exercise up in favor of another one that seems more promising.}
\par\vspace{6pt}\noindent
$\Rightarrow$ \rule{\linewidth}{0.4pt}

% ===================== Q4 =====================
\qno{4}\instr{윗글의 내용과 일치하지 \underline{않는} 것은?}
\begin{choices}
\item ① People often continue behaviors that produce some kind of benefit.
\item ② A child praised for doing an exercise correctly is likely to repeat that way.
\item ③ A dancer may abandon an exercise that seems to produce no result.
\item ④ People tend to act in ways that bring no reward and create discomfort.
\item ⑤ Consequences can influence how humans adapt their behavior.
\end{choices}

\end{document}
